%============================================================
% CH1 · PATTERN 2 — READ/WRITE POINTERS: THEORY
%============================================================
\subsection{Pattern Theory}
\label{ssec:ch1-p2-theory}

\begin{patternbox}{Pattern 2 at a Glance}
To \textbf{remove} elements from an array \textbf{in place} --- keep only those that pass a test, in their original order --- run two indices in the same direction: a \textbf{read pointer} $r$ that visits every element once, and a \textbf{write pointer} $w$ that marks where the next kept element goes. The array splits into three zones:

\begin{center}
\begin{tikzpicture}
    \dparray{z}{0}{}{3/akeep, 4/akeep, 9/akeep, $-1$/aout, $-5$/aout, 9/aout, 2/aread, $-6$/acell}
    \aptr{z3}{w}{accentorange}
    \aptr{z6}{r}{successgreen}
    \draw[thick, successgreen, decorate, decoration={brace, amplitude=4pt}] ([yshift=1.5mm]z0.north west) -- ([yshift=1.5mm]z2.north east) node[midway, above=4pt, font=\scriptsize\bfseries, text=successgreen] {kept $[0, w)$};
    \draw[thick, midgray, decorate, decoration={brace, amplitude=4pt}] ([yshift=1.5mm]z3.north west) -- ([yshift=1.5mm]z5.north east) node[midway, above=4pt, font=\scriptsize\bfseries, text=midgray] {free $[w, r)$};
    \draw[thick, primaryblue, decorate, decoration={brace, amplitude=4pt}] ([yshift=1.5mm]z6.north west) -- ([yshift=1.5mm]z7.north east) node[midway, above=4pt, font=\scriptsize\bfseries, text=primaryblue] {unread $[r, n)$};
    \node[note, anchor=west, text width=3.4cm, align=left] at (7.6,0) {input $3,-1,4,-1,-5,9,2,-6$, keep $x \ge 0$; the next step copies $2$ into the first free cell};
\end{tikzpicture}
\end{center}

\begin{recurrencebox}
\centering
$w \gets 0$\\[2pt]
for $r = 0, 1, \dots, n-1$:\quad if \textsc{keep}($a_r$, the written part $a[0..w)$) then $a_w \gets a_r$, \ $w \gets w + 1$\\[2pt]
the result is $a[0..w)$ \qquad $\bigO{n}$ time, $\bigO{1}$ extra memory, stable
\end{recurrencebox}
It works because $w \le r$ at every moment: the write position never overtakes the read position, so no element is overwritten before it has been read.
\end{patternbox}

%------------------------------------------------------------
\subsubsection*{From Part A to Pattern 2}

Lesson~\ref{ssec:ch1-inplace} named the danger of in-place work --- a write that destroys a value still needed --- and listed read/write pointers as one cure. This pattern develops that cure into a reusable tool. Where Pattern~1 \emph{reads} an array and summarises it, Pattern~2 \emph{rewrites} it: the ``state'' is now a whole prefix of the array, $a[0..w)$, which is the output under construction.

\begin{diagrambox}
\begin{tikzpicture}
    \begin{scope}
        \dparray{c}{0}{copy}{3/akeep, 4/akeep, 9/akeep, 2/akeep}
        \node[note, anchor=west, text width=4.0cm, align=left] at (3.9,0) {with extra memory: \texttt{out.push\_back(x)} for every kept $x$ --- $\bigO{n}$ extra};
    \end{scope}
    \begin{scope}[yshift=-1.2cm]
        \dparray{d}{0}{in place}{3/akeep, 4/akeep, 9/akeep, 2/akeep, $-5$/aout, 9/aout, 2/aout, $-6$/aout}
        \node[note, anchor=west, text width=3.3cm, align=left] at (7.5,0) {\texttt{a[w++] = x}: the output is the prefix of length $w = 4$; the tail holds leftovers};
    \end{scope}
\end{tikzpicture}
\figcap{The same filter (keep $x \ge 0$ from $3, -1, 4, -1, -5, 9, 2, -6$) with a second array and in place. The in-place tail holds old values that the caller must ignore (or overwrite, as Move Zeroes does).}
\end{diagrambox}

%------------------------------------------------------------
\subsubsection*{When to recognise it --- and when not}

Use Pattern 2 when
\begin{itemize}[leftmargin=*]
    \item the statement says ``remove \ldots{} \textbf{in place}'', ``return the new length $k$; the first $k$ elements must hold the result'', ``$\bigO{1}$ extra memory'';
    \item the kept elements must keep their \textbf{relative order} (stability) --- swaps from both ends would destroy it;
    \item the decision for an element depends on the element itself, or on the \textbf{last few kept} elements (deduplication, at most $k$ copies --- Problem~\ref{ssec:ch1-remove-duplicates-ii});
    \item the removed elements must be moved to the end rather than discarded (Problem~\ref{ssec:ch1-move-zeroes}: compact, then fill the tail);
    \item a later element can \textbf{cancel} an earlier kept one (backspace, ``a digit deletes the closest letter on its left''): the kept prefix behaves like a stack (Exercise~E2.3).
\end{itemize}

\textbf{Do NOT use it} (in this form) when
\begin{itemize}[leftmargin=*]
    \item the output is \textbf{longer} than the input (duplicating or inserting): $w$ would overtake $r$. Write from the \emph{back} instead (Exercise~E2.1);
    \item order does not matter and you only need to split into two groups: two pointers from both ends do fewer writes (partitioning, Chapter~\ref{ch:partitioning});
    \item the decision needs information about \textbf{all} elements (``remove values that occur more than twice'' in an \emph{unsorted} array): count first (Pattern~5), then compact with the counts.
\end{itemize}

%------------------------------------------------------------
\subsubsection*{The framework and why it is correct}

\begin{table}[H]
\centering\small
\renewcommand{\arraystretch}{1.25}
\begin{tabularx}{\linewidth}{@{}>{\bfseries}l X@{}}
\toprule
Maintained & $a[0..w)$ = the elements of $a[0..r)$ that were kept, in their original order \\
Step & read $x = a_r$; if kept, write it to $a_w$ and advance $w$; always advance $r$ \\
Initial & $w = r = 0$: nothing read, nothing kept \\
Answer & the new length $w$; the result is $a[0..w)$ \\
Complexity & $n$ reads, at most $n$ writes: $\bigO{n}$ time, $\bigO{1}$ extra memory \\
\bottomrule
\end{tabularx}
\end{table}

\begin{invariantbox}[Invariant of the read/write loop]
Before the step with index $r$: (1) $w \le r$; (2) $a[0..w)$ consists of the kept elements of the \emph{original} $a[0..r)$, in their original order; (3) $a[r..n)$ still holds the original values.

\textbf{Initialisation.} $w = r = 0$: all three statements are about empty ranges or the untouched array.

\textbf{Maintenance.} By (3), the cell $a_r$ still holds the original value $x$. If $x$ is not kept, nothing is written; $r$ grows, so (1)--(3) still hold. If $x$ is kept, we write $x$ to $a_w$. By (1), $w \le r$, so the written cell is either a ``free'' cell in $[w, r)$ (already read, no longer needed) or $a_r$ itself; in both cases no unread cell $a_{r+1..n-1}$ changes, so (3) holds for $r+1$. The kept list of $a[0..r+1)$ is the kept list of $a[0..r)$ followed by $x$, which is exactly $a[0..w+1)$ now: (2) holds. Both pointers grow by one, so $w \le r$ is preserved.

\textbf{Termination.} At $r = n$, (2) says $a[0..w)$ is the kept subsequence of the whole array. Every element was read once and written at most once. \hfill$\blacksquare$
\end{invariantbox}

The decision ``kept?'' may look at the \emph{written} part $a[0..w)$ --- that part is final and correct by (2). It must \textbf{not} rely on the original value of a cell in $[w, r)$: such a cell may already have been overwritten.

%------------------------------------------------------------
\subsubsection*{Three kinds of keep rules}

\textbf{(a) The element alone.} ``Remove every $x = \mathit{val}$'', ``keep the non-zero elements'', ``keep the even ones''. The test is \texttt{keep(a[r])}.

\textbf{(b) The last kept elements.} In a \emph{sorted} array, ``keep at most $k$ copies of each value'' means: keep $x$ unless the last $k$ kept elements all equal $x$. Since the array is sorted and the kept part is sorted too, this is the single comparison $x \ne a[w-k]$ (or $w < k$).

\begin{diagrambox}
\begin{tikzpicture}
    \dpindex{ki}{1.35}{}{0}{8}
    \dparray{k}{0.6}{\texttt{a}}{0/akeep, 0/akeep, 1/akeep, 1/akeep, 1/aout, 1/aout, 2/aread, 3/acell, 3/acell}
    \aptr{k2}{w-2}{tabulationblue}
    \aptr{k4}{w}{accentorange}
    \aptr{k6}{r}{successgreen}
    \node[note, anchor=west, text width=3.6cm, align=left] at (8.3,0.6) {$k = 2$: $a_r = 2 \ne a_{w-2} = 1$, so $2$ is kept and written to index $4$};
\end{tikzpicture}
\figcap{Compare the candidate with the cell $k$ places before the \emph{write} pointer. That cell is part of the final answer and is never overwritten again.}
\end{diagrambox}

\textbf{(c) Cancellation.} A special element removes the \emph{last kept} element instead of being kept: $w \gets w - 1$. The kept prefix is then a stack whose top is $a_{w-1}$. Still $w \le r$, still $\bigO{n}$.

\begin{diagrambox}
\begin{tikzpicture}[x=1cm, y=1cm, sc/.style={cell, minimum width=7mm, minimum height=6mm, font=\small}]
    \foreach \x/\st/\w [count=\k from 0] in {5/{5}/1, 3/{5,3}/2, 0/{5}/1, 7/{5,7}/2, 0/{5}/1, 0/{}/0, 8/{8}/1}{
        \node[font=\scriptsize\bfseries, text=deepblue] at (\k*1.55,0.6) {read $\x$};
        \ifnum\w>0 \foreach \v [count=\j from 0] in \st { \node[sc, fill=basefill] at (\k*1.55,-\j*0.62) {\v}; }\fi
        \node[font=\scriptsize, text=darkgray] at (\k*1.55,-1.55) {$w = \w$};
    }
\end{tikzpicture}
\figcap{Cancellation on $5, 3, 0, 7, 0, 0, 8$, where each $0$ deletes the last kept element: the kept prefix $a[0..w)$ after every step (top = first cell). The result is the single element $8$.}
\end{diagrambox}

%------------------------------------------------------------
\subsubsection*{Variations map}

\begin{diagrambox}
\begin{tikzpicture}[
    var/.style={rectangle, rounded corners=2pt, draw=patternborder, fill=patternbg,
                font=\scriptsize, align=center, text width=3.25cm, minimum height=1.2cm, inner sep=3pt}]
    \node[var, fill=answerfill, draw=accentorange, text width=5cm, font=\scriptsize\bfseries] (root) at (0,0)
        {Read/write compaction\\\texttt{if (keep) a[w++] = a[r]}};
    \node[var] (v1) at (-5.4,-2) {\textbf{remove by value}\\keep $x \ne \mathit{val}$, then fill the tail (Problem~\ref{ssec:ch1-move-zeroes})};
    \node[var] (v2) at (-1.8,-2) {\textbf{at most $k$ copies}\\sorted input: $x \ne a_{w-k}$ (Problem~\ref{ssec:ch1-remove-duplicates-ii})};
    \node[var] (v3) at (1.8,-2)  {\textbf{first occurrences}\\unsorted, small values: a seen-array (E2.2)};
    \node[var] (v4) at (5.4,-2)  {\textbf{cancellation}\\$w{-}{-}$: the kept part is a stack (E2.3)};
    \node[var] (v5) at (-3.6,-3.6) {\textbf{expansion}\\output longer: write from the back (E2.1)};
    \node[var] (v6) at (0,-3.6)    {\textbf{per row of a grid}\\gravity with walls: reset $w$ (E2.4)};
    \node[var] (v7) at (3.6,-3.6)  {\textbf{normalise separators}\\rule on $a_{w-1}$, trim the end (E2.5)};
    \node[draw=patternborder, dashed, rounded corners=4pt, fit=(v1)(v4)(v5)(v7), inner sep=5pt] (grp) {};
    \draw[flow] (root.south) -- (grp.north);
\end{tikzpicture}
\figcap{One loop, many problems. The bottom row is where most bugs are: the write pointer would overtake the read pointer, the segments must be handled separately, or the last written element must be undone.}
\end{diagrambox}

\textbf{Twists and how to handle them}
\begin{itemize}[leftmargin=*]
    \item \textbf{The tail must hold something specific} (zeros, the removed values): after compaction, fill $a[w..n)$; or swap instead of copy (Problem~\ref{ssec:ch1-move-zeroes}, Approach 3).
    \item \textbf{Several deletion rules at once}: one read pointer, one write pointer, a combined test --- never several passes of \texttt{erase}.
    \item \textbf{A second input} (a sorted list of positions to delete, Exercise~E2.6; or a second array to merge, Chapter~\ref{ch:two_pointers}): two read pointers and one write pointer.
    \item \textbf{A grid}: apply the pattern to every row or column independently; obstacles split a row into independent segments.
    \item \textbf{Output longer than input}: count the final length first, then fill from the right with $w$ starting at the end; now $w \ge r$ is the safety invariant.
\end{itemize}

\begin{diagrambox}
\begin{tikzpicture}
    \dpindex{bi}{1.35}{}{0}{7}
    \dparray{b}{0.6}{\texttt{a}}{1/acell, 0/acell, 2/acell, 3/aread, 0/aout, 0/akeep, 0/akeep, 4/akeep}
    \aptr{b3}{r}{successgreen}
    \aptr{b4}{w}{accentorange}
    \node[note, anchor=west, text width=3.9cm, align=left] at (7.5,0.6) {expanding (each 0 doubled, array truncated to length 8): read and write move \emph{leftwards}, $w > r$, so again no unread cell is overwritten};
\end{tikzpicture}
\figcap{The mirror image of the pattern on $1,0,2,3,0,4,5,0$: green cells are final ($0,0,4$), the grey cell is free; the next step writes $3$ at $w = 4$. Both pointers run from the right end towards the left. Exercise~E2.1 works this out.}
\end{diagrambox}

%------------------------------------------------------------
\subsubsection*{Code templates}

\begin{lstlisting}[style=cpp]
#include <bits/stdc++.h>
using namespace std;

// (1) keep the elements with keep(x) == true, in order; returns the new length
template <class Pred>
int compact(vector<int>& a, Pred keep) {
    int w = 0;
    for (int r = 0; r < (int)a.size(); r++)
        if (keep(a[r])) a[w++] = a[r];
    return w;
}

// (2) sorted array: keep at most k copies of every value
int atMostK(vector<int>& a, int k) {
    int w = 0;
    for (int r = 0; r < (int)a.size(); r++)
        if (w < k || a[r] != a[w - k]) a[w++] = a[r];   // compare with the WRITTEN part
    return w;
}

// (3) write pointer as a stack: every 0 deletes the last kept element
int cancel(vector<int>& a) {
    int w = 0;
    for (int r = 0; r < (int)a.size(); r++) {
        if (a[r] != 0) a[w++] = a[r];                   // push
        else if (w > 0) w--;                            // pop
    }
    return w;
}

void show(const vector<int>& a, int len) {
    cout << len << ":";
    for (int i = 0; i < len; i++) cout << " " << a[i];
    cout << "\n";
}

int main() {
    vector<int> a = {3, -1, 4, -1, -5, 9, 2, -6};
    show(a, compact(a, [](int x) { return x >= 0; }));     // 4: 3 4 9 2
    vector<int> b = {1, 1, 2, 3, 3, 3, 4}, c = b;
    show(b, atMostK(b, 1));                                // 4: 1 2 3 4
    show(c, atMostK(c, 2));                                // 6: 1 1 2 3 3 4
    vector<int> d = {5, 3, 0, 7, 0, 0, 8};
    show(d, cancel(d));                                    // 1: 8
}
\end{lstlisting}

\begin{insightbox}[The library version]
\texttt{std::remove\_if(a.begin(), a.end(), pred)} is template (1) with the \emph{opposite} predicate (it removes where \texttt{pred} is true) and returns an iterator to the new end; \texttt{std::unique} is template (2) with $k = 1$. Neither shrinks the vector: follow them with \texttt{a.erase(newEnd, a.end())} (the ``erase--remove idiom''). Templates (2) with $k > 1$ and (3) have no library equivalent.
\end{insightbox}

All three templates were compared with straightforward versions that build a new vector (a map of counts for (2), an explicit stack for (3)) on 50000 random arrays.

%------------------------------------------------------------
\subsubsection*{Common mistakes}

\begin{warningbox}[Mistake 1 --- erasing inside the loop]
\begin{lstlisting}[style=snippet]
for (int i = 0; i < (int)a.size(); i++)
    if (a[i] == val) a.erase(a.begin() + i);   // WRONG: O(n^2), and skips the element after
\end{lstlisting}
After an erase the next element slides into position $i$ and the \texttt{i++} jumps over it: $[2, 2]$ with $\mathit{val} = 2$ leaves one 2 behind. Even when fixed, every erase shifts the tail: $\bigO{n^2}$.
\end{warningbox}

\begin{warningbox}[Mistake 2 --- comparing with the read side]
\begin{lstlisting}[style=snippet]
if (r < 2 || a[r] != a[r - 2]) a[w++] = a[r];   // WRONG: a[r-2] may already be overwritten
if (w < 2 || a[r] != a[w - 2]) a[w++] = a[r];   // RIGHT: the written part is final
\end{lstlisting}
On $[1,1,1,2,2,3]$ the wrong version returns length 4 with $[1,1,2,3]$: by the time it reads the second 2, cell $a_2$ already holds the first 2, so the second 2 looks like a third copy and is dropped. Problem~\ref{ssec:ch1-remove-duplicates-ii} shows the trace.
\end{warningbox}

\begin{warningbox}[Mistake 3 --- forgetting what the caller sees]
LeetCode judges only $a[0..k)$ where $k$ is the returned length; other judges (and Move Zeroes) require the whole array. Return $w$, and if the statement requires it, fill or shrink the tail (\texttt{a.resize(w)}).
\end{warningbox}

\begin{warningbox}[Mistake 4 --- expanding with a forward write pointer]
Duplicating elements left to right makes $w$ run ahead of $r$ and overwrite unread values: duplicating zeros in $[1, 0, 2, 3]$ forward yields $[1, 0, 0, 0]$ instead of $[1, 0, 0, 2]$. Go right to left (Exercise~E2.1).
\end{warningbox}

\begin{checklistbox}[``I'm stuck'' checklist for Pattern 2]
\begin{itemize}[leftmargin=*]
    \item[$\square$] Is the output a subsequence of the input (same order, some removed)? Then read/write pointers.
    \item[$\square$] What does the keep rule look at: the element, the last kept elements ($a_{w-1}, a_{w-k}$), or the original neighbours?
    \item[$\square$] Is $w \le r$ guaranteed? If elements are inserted, run both pointers from the back.
    \item[$\square$] What must the tail $a[w..n)$ contain afterwards? What must be returned?
    \item[$\square$] Do I need information about the whole array first (counts)? Then a first pass, then compaction.
    \item[$\square$] Is stability required? Copying is stable; swapping from both ends is not.
\end{itemize}
\end{checklistbox}
